\documentclass{article}

\usepackage[utf8]{inputenc}     % pdfLaTeX
\usepackage[T1]{fontenc}
\usepackage[magyar]{babel}
\usepackage{graphicx}

% --- Section formatting ---
\usepackage{titlesec}

% Allow deep section numbering and TOC
\setcounter{secnumdepth}{5}
\setcounter{tocdepth}{5}

\makeatletter

% ------------------------
% LEVEL 4: subsubsubsection
% ------------------------
\titleclass{\subsubsubsection}{straight}[\subsubsection]
\newcounter{subsubsubsection}[subsubsection]
\renewcommand\thesubsubsubsection{\thesubsubsection.\arabic{subsubsubsection}}

\renewcommand\subsubsubsection{%
  \@startsection{subsubsubsection}{4}{\z@}%
  {3.25ex \@plus1ex \@minus.2ex}%
  {1.5ex \@plus .2ex}%
  {\normalfont\normalsize\bfseries}%
}

\titleformat{\subsubsubsection}
  {\normalfont\normalsize\bfseries}
  {\thesubsubsubsection}{1em}{}

\titlespacing*{\subsubsubsection}{0pt}{3ex}{1ex}

% TOC entry for level 4
\newcommand*\l@subsubsubsection{\@dottedtocline{4}{7em}{4em}}

% --------------------------
% LEVEL 5: subsubsubsubsection
% --------------------------
\titleclass{\subsubsubsubsection}{straight}[\subsubsubsection]
\newcounter{subsubsubsubsection}[subsubsubsection]
\renewcommand\thesubsubsubsubsection{\thesubsubsubsection.\arabic{subsubsubsubsection}}

\renewcommand\subsubsubsubsection{%
  \@startsection{subsubsubsubsection}{5}{\z@}%
  {3.25ex \@plus1ex \@minus.2ex}%
  {1.5ex \@plus .2ex}%
  {\normalfont\normalsize\bfseries}%
}

\titleformat{\subsubsubsubsection}
  {\normalfont\normalsize\bfseries}
  {\thesubsubsubsubsection}{1em}{}

\titlespacing*{\subsubsubsubsection}{0pt}{3ex}{1ex}

% TOC entry for level 5
% Use level 4 internally for safety (LaTeX only knows TOC levels 1-4)
\newcommand*\l@subsubsubsubsection{\@dottedtocline{4}{10em}{5em}}

\makeatother

\title{Vallások}
\author{Haydu Lőrinc, Lengyel Zoltán Péter, Rovenszky Robin}
\date{Legutóbb frissítve: 2026. Március 16.}

\begin{document}

\maketitle
\tableofcontents
\newpage

\section{Iszlám vallás}
\begin{itemize}

    \item 'iszlám' - alávetés / megadás Allahnak
    \item Iszlám - vallás
    \item Muszlim - követők
    
    \item Szent könyve: Korán
    \begin{itemize}
    
        \item Allah az egyetlen Isten, mindenki az ő prófétája
        \item Napi ötszöri ima Mekka felé, általában rituális tisztálkodás/fürdő előzi meg
        \item Alamizsna a szegényeknek
        
        \item Böjt: Ramadan,
        \begin{itemize}
            \item 60 napig tart,
            \item Csak naplemente után étkeznek.
        \end{itemize}
        
        \item Mindenkinek egyszer el kell zarándokolnia Mekkába megnézni a Fekete követ.
        
    \end{itemize}
    
    \item További előírások 
    \begin{itemize}
        \item Sertéshús fogyasztása tilos
        \item Többnejűség korlátozása
        \item Szeszes ital fogyasztása tilos
        \item Szerencsejáték és tisztátlan húsok fogyasztása
    \end{itemize}
    
    \item Vallási szerkezet 
    \begin{itemize}
        \item Vallás elején Allah 
        \item Alatta próféták 
        
        \item alatta kalifa
            \begin{itemize}
                \item Vallási vezető
                \item Állami vezető
                \item Mohamed leszármazottjának tartják
            \end{itemize}
            
        \item defterdár - pénzügyekkel foglalkozik
    \end{itemize}
    
    \item Dzsihád - Szent Háború
    \begin{itemize}
        \item Hitterjesztés alapvetően békés módon
        \item De fegyveresen is
    \end{itemize}
    
\end{itemize}

\subsection{Arabok}

\begin{itemize}
    \item Törzsi szervezet
    \item Nomád életforma
    \item Sokistenhit
    \item Kereskedelem (Tömjénút)
    
    \begin{itemize}
        \item Városok (pl. Mekka $\rightarrow$ szent város)
    \end{itemize}
    
\end{itemize}

\subsection{Mohamed próféta (kb. 570-632)}

\begin{itemize}
    \item Próféta - követ/küldönc
    \item Monoteizmust hirdeti - Allah a legfőbb isten
    \item Fellép a kapzsiság ellen, elítéli az uzsorát és a harácsoló életmódot
    \item Támogatja a szegényeket
    \item 622 = hidzsra ('kivonulás') $\rightarrow$ iszlám időszámítás kezdete $\rightarrow$ Mohamedet elűzték Mekkából
\end{itemize}

\subsection{Iszlám kultúra}

\begin{itemize}
    
    \item Magas színvonalú műveltség
    \begin{itemize}
        \item Természettudományok
        \item Csillagászat
        \item Matematika
        \item Könyvtárak
        \item Anatómia
    \end{itemize}
    
    \item Elterjedése: 
    \begin{itemize}
        \item Törökország
        \item Spanyolország
        \item Észak-Afrika
        \item Közel-Kelet
        \item Dél-nyugat-Ázsia
    \end{itemize}
    
\end{itemize}

\subsection{Iszlám történelem és művészet korszakolása}

\begin{itemize}
    \item Mohamed 570-632
    \item Négy igaz úton járó kalifa 632-661
    \item Omajjád Kalifátus 661-750
    \item Abbászida Kalifátus 749-1259
    \item Mamelukok 1250-1722
    \item Oszmán Birodalom 1517-1923
\end{itemize}

\subsection{Iszlám építészet}

\begin{itemize}

    \item Kezdetben a kalifátusok nem a hatalmukat akarták kifejezni, hanem a funkció volt a jelentős
    \begin{itemize}
        \item Óváros 
        \item Hammám - fürdő
        \begin{itemize}
            \item Férfi és női napok (meztelenség)
            \item Középen általában nyolcszögletű fürdő
            \item kupolás szerkezetek támpillérekkel
        \end{itemize}
        \item Szeráj -  uralkodói palota \begin{itemize}
            \item Hárem - női lakrész
            \item lakószobák - 
            \item tornácok
            \item fürdőhelyiségek
            \item külső és belső kert
            \item először sivatagi paloták, majd luxuskomplexumok
            \item Példák: \begin{itemize}
                \item Alhambra
            \end{itemize}
        \end{itemize}
        \item Karavánszeráj - kereskedők szállása
        
        \item Mecset - imádság helye, nem szent hely
        \begin{itemize}
            \item Imaszőnyegek
            \item Imafülke
            \item Tájolófal (merre van Mekka)
            
            \item Példák:
            \begin{itemize}
                \item Mekkai nagymecset (Masjid al-Haram)
                \item Hagia Szofia (Ayasofya)
                \item Szulejmán szultán mecset
                \item Ahmed szultán mecsete
            \end{itemize}
        \end{itemize}
        
        \item Dzsámi - imádság helye, nem szent hely (most 200\%-al nagyobb)
        \begin{itemize}
            \item Péntekenként használták
            \item Imafülke
            \item Tájolófal
        \end{itemize}
        
        \item Medresze - iskola 
        \begin{itemize}
            \item Felsőoktatási intézmény és Általános iskola is lehet
            \item Teológia tudományának iskolája
        \end{itemize}
        
        \item Ispotály - kórház
            \item Kümbet és türbe 
                \begin{itemize}
                    \item Sírhely
                    \item Türbe: szarkofág köré épített
                    
                    \item Kümbet: nagyobb síremlék 
                    \begin{itemize}
                        \item pl. Taj Mahal
                    \end{itemize}
            \end{itemize}
    \end{itemize}
    
    \item Cordova
    \item Bagdad
    \item Damaszkusz
    \item Jeruzsálem
    \item Isztambul (Bizánc)
    
\end{itemize}

% missing table explaining jews christians and muslims from somewhere in the note

\section{Keresztények}
\subsection{Szerzetesrendek}
\subsubsection{Bencések}
\subsubsubsection{Pannonhalma}
\begin{itemize}
    \item 996. Szent Benedek követői alapították a monostort
    \begin{itemize}
        \item Szent Asztrik (megalapítja Pannonhalmát, neveli Szent Istvánt)
    \end{itemize}
    \item 15. században épült kerengő
    \begin{itemize}
        \item Szent Benedek Kápolna
    \end{itemize}
    \item Építészeti jellemzői
    \begin{itemize}
        \item Román stílus
        \begin{itemize}
            \item Lőrésszerű ablakok
            \item vastag falak (pl Jáki templom, Városligetben másolata)
            \item bélletes kapu (befelé szűkülő kapu)
            \item 1 főhajó + 1 mellékhajó
        \end{itemize}
        \item Korai gótika 
        \begin{itemize}
            \item Csúcsíves
            \item Támpillérek
            \item Rózsaablakok
            \item Környezeténél magasabb
            \item keskeny falak
            \item finom vonások
            \item pl Notre Dame
        \end{itemize}
        \item Késő gótika
        \begin{itemize}
            \item kevésbé magas
            \item vízköpő (díszelem a külső részeken)
        \end{itemize}
    \end{itemize}
    \item Apátság alapítólevele 1045.
\end{itemize}

\subsubsection{Ciszterciek}
\begin{itemize}
    \item Citeaux (Francia tartomány)
     \item 1098. Szent Róbert
     \item 1102. Szent Bernát - virágoztatja fel
     \item Kezdetben fehér $\rightarrow$ barna $\rightarrow$ fekete viselet (csuha)
     \item Carta Cartiatis (szeretet okirata) rend nűködésének regulája
     \item pl. Gödöllő
\end{itemize}

\subsubsection{Egyéb szerzetesrendek}
\begin{itemize}
    \item Premontreiek
    \item Karthauziak
    \item Ferencesek
    \item Domonkosok
    \item Johanniták
\end{itemize}


\end{document}